\section{Gamma-free curvature, logarithms, and a probabilistic transfer}
\label{sec:crossfield-transfers}

The pure-exponential base change is stable under every full-rank linear
projection.  This produces explicit gamma-free families among classical
logarithmic special functions and transfers the same algebraic matroid to a
probability law.

\subsection{A geometric-grid curvature family}

Fix an integer \(q\ge2\).  For \(n\ge0\), define
\begin{align*}
 A_n&=E_1(q^{n+2})-2E_1(q^{n+1})+E_1(q^n),\\
 B_n&=\Ei(q^{n+2})-2\Ei(q^{n+1})+\Ei(q^n),\\
 C_n&=\Ci(q^{n+2})-2\Ci(q^{n+1})+\Ci(q^n),\\
 S_n&=\Si(q^{n+2})-2\Si(q^{n+1})+\Si(q^n).
\end{align*}

\begin{theorem}[Gamma-free geometric curvature]
\label{thm:geometric-curvature}
The countable family
\begin{equation}\label{eq:geometric-curvature-family}
 \{A_n,B_n,C_n,S_n:n\ge0\}
\end{equation}
is algebraically independent over \(\Eexp\), in the finite-subset sense.
\end{theorem}

\begin{proof}
For positive \(x\), the connection formulas are
\[
 \Ein(x)=\gamma+\log x+E_1(x),\qquad
 \Ein(-x)=\gamma+\log x-\Ei(x),
\]
and
\[
 \frac{\Ein(ix)+\Ein(-ix)}2=\gamma+\log x-\Ci(x),
 \qquad
 \frac{\Ein(ix)-\Ein(-ix)}{2i}=\Si(x).
\]
The second-difference row \((1,-2,1)\) kills both the constant \(\gamma\)
and the affine sequence \(\log(q^n)=n\log q\).  Hence every displayed
quantity is a linear form in \(\Ein\)-values on one of the four disjoint
orbits \(\{q^j\}\), \(\{-q^j\}\), \(\{iq^j\}\), and \(\{-iq^j\}\).
Within each block, every finite family of second-difference rows is linearly
independent: inspect the largest terminal index in a linear combination.
The four blocks have disjoint support.  Lemma~\ref{lem:linear-image-ein}
therefore proves joint algebraic independence over \(\Eexp\).
\end{proof}

The geometric grid is essential to the uncorrected formula: on a general
grid one must add the logarithmic second difference
\(\log(x_{n+2}x_n/x_{n+1}^2)\).  Omitting it would reintroduce the connection
term and invalidate the linear-image argument.

\subsection{Logarithms of rational period expressions}

\begin{proposition}[Global logarithm principle]
\label{prop:global-log-principle}
Let \(\lambda_1,\ldots,\lambda_s\in\Qbar^\times\) be distinct, and let
\(R\in\Eexp(X_1,\ldots,X_s)\) be nonconstant.  Suppose
\[
 r=R(\Ein(\lambda_1),\ldots,\Ein(\lambda_s))
\]
is defined and nonzero.  Every solution of \(e^w=r\), hence every branch of
\(\log r\), is transcendental.
\end{proposition}

\begin{proof}
If \(w\) were algebraic, then \(e^w\in\Eexp\).  On the other hand,
Theorem~\ref{thm:pure-exp-base} makes the \(\Ein(\lambda_j)\) independent
indeterminates over \(\Eexp\), so a nonconstant rational expression in
them cannot belong to \(\Eexp\).  This contradiction proves the claim.
\end{proof}

The proposition applies after every invertible Fourier transform in
\eqref{eq:fourth-root-transform}.  For example, every logarithm of each of
\(\Si(a),\Cin(a),\Shi(a),\Cinh(a)\), with
\(a\in\Qbar^\times\), is transcendental.  It also proves the transcendence
of every branch of \(\log\Ein(1)\) without an exceptional-value analysis;
the non-ultra-Liouville refinement in \cref{sec:log} still uses the
quantitative theorem of Fischler--Rivoal.

More generally, let
\(G(W,e^{\beta_1W},\ldots,e^{\beta_mW})\) be a rational expression over
\(\Qbar\), where the \(\beta_j\) are algebraic.  Every solution \(w\) of
\[
 G(w,e^{\beta_1w},\ldots,e^{\beta_mw})=r
\]
is transcendental whenever both sides are defined: an algebraic \(w\)
would make the left side an element of \(\Eexp\).  Logarithms and Lambert
\(W\)-preimages of \(r\) are elementary special cases.

\subsection{The Dickman log-Laplace matroid}

Let \(\Phi\) denote the Laplace transform of the normalized Dickman law.
The classical de Bruijn formula is
\begin{equation}\label{eq:dickman-laplace}
 \Phi(s)=\exp\left(\int_0^1\frac{e^{-sx}-1}{x}\,dx\right)
        =e^{-\Ein(s)}\qquad(s>0);
\end{equation}
see, for example, \cite{Moree2014}.  Thus the L\'evy exponent, rather than
the transform value itself, is an \(\Ein\)-period.

\begin{corollary}[Dickman log-Laplace periods]
\label{cor:dickman-matroid}
For distinct positive algebraic \(s_1,\ldots,s_N\), the numbers
\[
 -\log\Phi(s_j)=\Ein(s_j)\qquad(1\le j\le N)
\]
are algebraically independent over \(\Eexp\).  If these parameters label
the vertices of an oriented graph, then the log-ratios
\[
 \log\frac{\Phi(s_{t(e)})}{\Phi(s_{s(e)})}
 =-\Ein(s_{t(e)})+\Ein(s_{s(e)})
\]
have precisely the signed cycle relations and no others.
\end{corollary}

\begin{proof}
Use \eqref{eq:dickman-laplace}, Theorem~\ref{thm:pure-exp-base}, and the
incidence-matrix argument of Theorem~\ref{thm:graphic-period}.
\end{proof}

This is an exact transfer to the logarithmic Laplace coordinates of a
classical smooth-number distribution.  It gives no arithmetic conclusion
about the values \(\Phi(s_j)=e^{-\Ein(s_j)}\) themselves, which lie outside
the finite-point Siegel--Shidlovskii mechanism.
